\documentclass[11pt,a4paper]{article}

% ── Packages ──────────────────────────────────────────────
\usepackage[margin=2.2cm]{geometry}
\usepackage{fontspec}
\usepackage{xeCJK}
\setCJKmainfont{Noto Serif CJK SC}
\setCJKsansfont{Noto Sans CJK SC}
\setCJKmonofont{Noto Sans Mono CJK SC}
\usepackage{graphicx}
\usepackage{listings}
\usepackage{xcolor}
\usepackage{booktabs}
\usepackage{enumitem}
\usepackage{hyperref}
\usepackage{fancyhdr}
\usepackage{titlesec}
\usepackage{tabularx}

% ── Style ─────────────────────────────────────────────────
\definecolor{codebg}{HTML}{F5F5F5}
\definecolor{codeblue}{HTML}{2563EB}
\definecolor{codegray}{HTML}{6B7280}
\definecolor{accent}{HTML}{1E40AF}

\lstset{
  basicstyle=\small\ttfamily,
  backgroundcolor=\color{codebg},
  keywordstyle=\color{codeblue}\bfseries,
  commentstyle=\color{codegray}\itshape,
  stringstyle=\color{red!60!black},
  breaklines=true,
  frame=single,
  framerule=0.3pt,
  rulecolor=\color{gray!40},
  xleftmargin=2em,
  framexleftmargin=1.5em,
  numbers=left,
  numberstyle=\tiny\color{codegray},
  numbersep=6pt,
  tabsize=4,
  language=C,
  morekeywords={ToolDef,ToolResult,Agent,Context,ContextPolicy,LLMResponse,
    cJSON,MessageList,SessionMeta,MemoryType,HookEvent,ToolFn,ToolCallView},
}

\hypersetup{
  colorlinks=true,
  linkcolor=accent,
  urlcolor=accent,
}

\titleformat{\section}
  {\Large\bfseries\color{accent}}{\thesection}{1em}{}
\titleformat{\subsection}
  {\large\bfseries\color{accent!80}}{\thesubsection}{1em}{}

\pagestyle{fancy}
\fancyhf{}
\fancyhead[L]{\small\color{codegray}CS2313 Project 2 Report}
\fancyhead[R]{\small\color{codegray}C-Agent: A Terminal AI Coding Agent}
\fancyfoot[C]{\thepage}
\renewcommand{\headrulewidth}{0.4pt}

% ── Document ──────────────────────────────────────────────
\begin{document}

\begin{center}
  {\LARGE\bfseries\color{accent} C-Agent: A Terminal AI Coding Agent}\\[6pt]
  % {\Large C-Agent: A Terminal AI Coding Agent}\\[10pt]
  {\large\color{codegray} CS2313 --- Operating Systems Course Project}\\[6pt]
  {\normalsize Hengtao Wu \quad 524031910038 \quad \href{https://github.com/EternalWavee/c-agent}{\texttt{GitHub}}}\\[3pt]
  {\scriptsize Source code is available at \href{https://github.com/EternalWavee/c-agent}{\texttt{github.com/EternalWavee/c-agent}}}\\[4pt]
  {\normalsize June 2026}
\end{center}
\vspace{-1em}

% ──────────────────────────────────────────────────────────
\section{Architecture Overview}
% ──────────────────────────────────────────────────────────

C-Agent is a pure-C terminal coding assistant that follows the \textbf{ReAct}
(Reason + Act) loop pattern. The user types a natural-language request; the
agent forwards it to a remote LLM, executes whatever tools the LLM requests,
feeds the results back, and prints the final answer. Although the handout only
requires a basic coding agent, I intentionally treated the project as a small
operating-system runtime for AI agents: it schedules tool calls, isolates file
access, persists sessions, reclaims context memory, and exposes an extensible
interface for future capabilities.

\subsection{Module Map}

\begin{center}
\small
\begin{tabularx}{\textwidth}{l l X}
\toprule
\textbf{Module} & \textbf{Key Files} & \textbf{Responsibility} \\
\midrule
Agent      & \texttt{agent/agent.c}      & ReAct loop, context reclaim, UI sync \\
LLM Client & \texttt{agent/llm\_client.c} & HTTP request construction, JSON response parsing \\
Context    & \texttt{context/}           & Token budget tracking, offload \& summary policies \\
Executor   & \texttt{tools/executor.c}   & Serial / parallel tool dispatch \\
Registry   & \texttt{tools/registry.c}   & ToolDef dispatch table \\
Sandbox    & \texttt{tools/sandbox.c}    & Path containment via \texttt{realpath} \\
Session    & \texttt{session.c}          & Checkpoint + append-only log persistence \\
Memory     & \texttt{memory.c}           & 6-type typed memory, proactive remember, manual consolidation \\
Skills     & \texttt{skills.c}           & Progressive skill discovery and loading \\
UI         & \texttt{ui/}               & Render thread, spinner, markdown output \\
Hooks      & \texttt{hooks.c}           & Tool lifecycle event bus \\
\bottomrule
\end{tabularx}
\end{center}

\subsection{Data Flow}

A single user request traverses the system as follows:

\begin{enumerate}[leftmargin=2em]
  \item \texttt{main.c} reads input; \texttt{cmd\_dispatch} handles \texttt{/} commands.
  \item \texttt{agent\_chat} pushes the user message into the \texttt{Context}.
  \item \texttt{ctx\_reclaim} checks token budget; triggers offload/summary if needed.
  \item \texttt{llm\_chat} builds an HTTP/1.1 POST to \texttt{/api/v1/chat/completions},
        including all registered \texttt{ToolDef} schemas discovered from the registry.
  \item While the LLM returns tool calls: \texttt{executor\_run\_tools} dispatches them
        (parallel if all read-only), pushes results, reclaims context, and calls the LLM again.
  \item The final text reply passes through \texttt{ui\_idle()} (barrier) then
        \texttt{ui\_print\_markdown()}.
\end{enumerate}



% ──────────────────────────────────────────────────────────
\section{Key Design Decisions}
% ──────────────────────────────────────────────────────────

\subsection{Dispatch Table for Tool Registry}

Inspired by OS syscall tables and VFS operation tables, each tool is a
self-contained \texttt{ToolDef} struct bundling name, description, JSON schema,
execution function pointer, and a \texttt{read\_only} scheduling hint. The
registry is a flat array populated once at boot; at runtime, tool lookup is a
linear scan over at most 16 entries. This design means \emph{adding a new tool
requires zero changes to the agent loop or LLM client}---a critical property
as the project grew from 1 tool (Part 1) to 16 tools (Part 3 + extensions).

\begin{lstlisting}[title={ToolDef interface (tools/tools.h)}]
typedef ToolResult (*ToolFn)(cJSON *args);

typedef struct {
    const char *name;
    const char *desc;
    const char *param_schema;  // JSON Schema string
    ToolFn exec;
    bool read_only;            // scheduling metadata
} ToolDef;
\end{lstlisting}

\subsection{Conservative Parallel Scheduling}

The executor uses a simple but correct rule: if \emph{all} tools in a batch are
registered and \texttt{read\_only}, dispatch them as concurrent pthreads;
otherwise fall back to serial execution. Each task writes only to its own slot
\texttt{tasks[i].result}; the main thread reads all slots only after every
worker joins. This preserves the \textbf{request-order invariant}
(\texttt{out\_msgs[i]} corresponds to \texttt{tool\_calls[i]}) regardless of
completion order---analogous to how CPUs may execute out-of-order but retire
in-order.

\subsection{Two-Layer Context Reclamation}

Context management uses a policy engine with pluggable strategies:

\begin{itemize}[leftmargin=2em]
  \item \textbf{Offload policy} (lossless): When token usage exceeds
        \texttt{OFFLOAD\_THRESHOLD} (default 0.8), large tool outputs are
        written to session-local files under \texttt{.agent/offload/} and
        replaced with compact placeholders containing recovery hints. The
        policy also maintains an offload manifest recording each artifact's
        path, size, and source message, so the agent can inspect what was
        paged out and call \texttt{read\_file} to retrieve the full content on
        demand.
  \item \textbf{Summary policy} (lossy): When usage exceeds
        \texttt{SUMMARY\_THRESHOLD}, a dedicated LLM call compresses old
        messages into a structured handoff with sections like
        \texttt{## Goal}, \texttt{## Completed}, \texttt{## Next Steps}.
\end{itemize}

Offload runs first because it is lossless and cheap; summary handles the
remaining verbosity. This staged design mirrors virtual memory's
\emph{page-out before OOM-kill} philosophy.

\subsection{Sandbox: \texttt{realpath} over String Checks}

Path validation uses \texttt{realpath(3)} rather than string-level
\texttt{".."} detection. This correctly handles symlink attacks, paths like
\texttt{foo/../../etc/passwd}, and absolute paths outside the workspace. For
write targets that do not yet exist, the parent directory is resolved and
checked, then the leaf is grafted back by string concatenation.

\subsection{Session Persistence: Checkpoint + Append-Only Log}

Sessions are organized as directories under
\texttt{.agent/sessions/<session\_id>/}. Each directory contains metadata for
UI display, a \texttt{checkpoint.json} full snapshot, and an append-only
\texttt{log.jsonl} delta stream. Recovery replays the log on top of the latest
checkpoint, while automatic checkpointing every 10 messages bounds recovery
time. This directory-level organization also gives context offload a
session-local namespace, so offloaded artifacts from different conversations do
not get mixed together. All writes are atomic (write to \texttt{.tmp}, then
\texttt{rename}) to survive crashes mid-write.


% ──────────────────────────────────────────────────────────
\section{Extensions Beyond the Handout}
% ──────────────────────────────────────────────────────────

The extra features were not added simply to increase the tool count. They were
chosen to make the agent usable as a long-running development assistant. In
particular, session persistence makes the REPL recoverable, memory preserves
stable facts across conversations, skills reduce system-prompt bloat, web tools
let the agent inspect current documentation, hooks decouple cross-cutting
behavior from the executor, and markdown rendering makes the terminal output
readable enough for daily use.

At a higher level, these extensions turn the project from a one-shot ReAct
demo into a \textbf{stateful agent runtime}. The design adds five production
properties that are not usually visible in minimal agents: durable identity
(memory), durable execution history (sessions), bounded context growth
(offload/summary), capability modularity (skills and tool registry), and
operational observability (hooks and synchronized UI). This framing was useful
because each feature supports a system-level invariant rather than existing as
an isolated bonus command.

\begin{center}
\small
\begin{tabularx}{\textwidth}{l X}
\toprule
\textbf{Extension} & \textbf{Design Summary} \\
\midrule
SubAgent &
  Thread pool (max 8). Each child is a \texttt{silent} Agent with its own
  Context. \texttt{subagent\_wait} joins the thread and returns the result.
  Intermediate output stays out of the parent's context window, so delegation
  does not pollute the parent conversation. \\[4pt]
Session &
  Append-only JSONL log with periodic checkpoints. Interactive TUI selector
  with arrow-key navigation, timestamps, and message counts. Supports
  rename, tag, delete, and restore. Old single-file format auto-migrated,
  making the storage layer backward compatible. \\[4pt]
Memory &
  6 typed files under \texttt{.agent/memory/}. The prompt asks the agent to
  proactively remember durable user facts, preferences, project workflows, and
  explicit corrections. \texttt{/memory} performs manual consolidation. On
  startup, the agent injects only selected fact/preference memories into the
  system prompt, so stable user identity and preferences are available in a new
  session without loading the entire memory store. \\[4pt]
Skill &
  Progressive disclosure: startup injects only skill names and descriptions.
  \texttt{load\_skill} fetches the full \texttt{SKILL.md} on demand. Project
  skills override user skills with the same name, giving the runtime a
  namespace and override model similar to configuration layers. \\[4pt]
Web Access &
  \texttt{web\_fetch} (curl wrapper) and \texttt{web\_search} (Bing RSS).
  For JavaScript-heavy sites, the installed \texttt{agent-browser} skill can
  render pages through Chrome. Private and localhost URLs are blocked, keeping
  web access useful without exposing local services. \\[4pt]
Hooks &
  Event bus with \texttt{HOOK\_BEFORE\_TOOL} and \texttt{HOOK\_AFTER\_TOOL}.
  Memory observation is registered as a hook, decoupling it from the
  executor and making future cross-cutting features easier to add. \\[4pt]
Dangerous Command Filter &
  \texttt{bash} tool checks 16 patterns (\texttt{rm -rf /}, fork bomb,
  \texttt{curl|sh}, etc.) before execution. \\[4pt]
Markdown UI\footnotemark &
  The UI uses MD4C to render markdown in the terminal and keeps spinner/tool
  status output synchronized through a dedicated render thread. \\
\bottomrule
\end{tabularx}
\end{center}
\footnotetext{\footnotesize Thanks to \href{https://github.com/DaveGamble/cJSON}{\texttt{cJSON}}, the lightweight C JSON parser used in this project, and \href{https://github.com/mity/md4c}{\texttt{MD4C}}, the markdown parser used for terminal rendering.}

% ──────────────────────────────────────────────────────────
\section{Challenges and Solutions}
% ──────────────────────────────────────────────────────────

\textbf{Concurrent tool output interleaving.}
When multiple read-only tools complete simultaneously, their
\texttt{ui\_tool\_done} calls could interleave terminal output. The solution:
a dedicated render thread owns all stdout via a mutex-protected event queue;
the main thread only posts events. \texttt{ui\_idle()} acts as a barrier
before the main thread reclaims output authority.

\textbf{Context budget explosion.}
Early versions triggered summary on every turn after the threshold was
crossed, causing an LLM call per turn and sometimes hitting model rate limits.
Another bug appeared when a recent browser snapshot was too large: summary kept
recent messages, so the context could remain above the 32K window forever. Fix:
\texttt{policy\_offload} now scans all tool messages, including recent large
outputs, and \texttt{ctx\_reclaim} stops after the first policy that actually
reduces token count.

\textbf{Offload recovery loops.}
If the LLM sees an offload placeholder and calls \texttt{read\_file} to
recover it, the recovered content could be offloaded again next turn. Fix:
\texttt{is\_offload\_recovery\_read} detects when a \texttt{read\_file} targets
an offload path and skips re-offloading that message.

\textbf{UTF-8 in tool output.}
Some shell commands produce invalid UTF-8 (e.g.\ \texttt{xxd}, binary files).
This corrupted the JSON body sent to the LLM. Fix:
\texttt{sanitize\_utf8\_for\_json} in \texttt{llm\_client.c} replaces
invalid sequences with \texttt{?} before serialization.

\textbf{Session migration and selector usability.}
The session format evolved from a single JSON file to a directory-based
checkpoint+log format. \texttt{migrate\_old\_sessions} detects legacy files,
converts them to the new layout, and removes the originals transparently on
startup. The interactive selector also required practical fixes: empty
\texttt{0 msg} sessions are filtered, session sizes are shown, and Chinese
session names are aligned by UTF-8 display width instead of raw byte length.

\textbf{Search and browser automation loops.}
A naive Bing HTML parser failed on current result pages, causing the LLM to
repeat \texttt{web\_search}, \texttt{web\_fetch}, and offload recovery in a
loop. The final version uses Bing RSS (\texttt{format=rss}) for structured
\texttt{title/link/description} entries and explicitly tells the agent not to
fetch Bing result pages repeatedly. For SPA pages, the installed
\texttt{agent-browser} skill renders the page with Chrome; on Linux it must run
with \texttt{--no-sandbox}, which is documented as an operational constraint.

% ──────────────────────────────────────────────────────────
\section{Effort Estimation vs.\ Reality}
% ──────────────────────────────────────────────────────────

\begin{center}
\small
\begin{tabular}{lrrr}
\toprule
\textbf{Phase} & \textbf{Estimated} & \textbf{Actual} & \textbf{Notes} \\
\midrule
Part 1 Phase A (LLM + bash)     & 4h  & 6h  & JSON parsing edge cases \\
Part 1 Phase B (ReAct loop)     & 2h  & 2h  & Straightforward extension \\
Part 1 Phase C (REPL + TUI)     & 3h  & 4h  & Render thread sync tricky \\
Part 2 Phase A (Registry)       & 3h  & 3h  & Clean refactor \\
Part 2 Phase B (Parallel)       & 2h  & 3h  & TSan debugging \\
Part 3 Context Management       & 4h  & 7h  & Recent huge tool outputs, recovery loops \\
Extensions (all)                & 8h  & 18h & Memory, sessions, skills, browser/web, markdown UI \\
\midrule
\textbf{Total}                  & 26h & 45h & \\
\bottomrule
\end{tabular}
\end{center}

The largest surprise was the \textbf{interaction surface} between extensions.
Session persistence needed to coordinate with context offload (which session
directory to use), memory needed hooks in the executor, web/browser outputs
changed the behavior of context reclamation, and the system prompt had to
dynamically assemble memory, skills, web guidance, and context hints. Each
extension was individually simple, but their composition required careful
ordering and shared-state management. In retrospect, the project was less about
writing isolated C functions and more about designing invariants: every tool
result must remain recoverable, every persisted session must be replayable, and
every background subsystem must not corrupt the main conversation.

% ──────────────────────────────────────────────────────────
\section{Course Improvement Suggestions}
% ──────────────────────────────────────────────────────────

\begin{enumerate}[leftmargin=2em]
  \item \textbf{Provide a proxy-aware test harness.} Students behind HTTP
        proxies (common in lab environments) will see mysterious 502 errors.
        The test runner should unset proxy variables or set
        \texttt{NO\_PROXY=localhost} automatically.

  \item \textbf{Add a concurrency stress test.} The current parallel test
        checks correctness and coarse speedup, but does not exercise
        thread-safety under contention. A test that fires 8 concurrent
        \texttt{write\_file} calls to the same file would catch missing locks.

  \item \textbf{Clarify the \texttt{tools\_advertised} test expectation.}
        The test checks for an exact 4-tool set, but extensions naturally
        register more tools. A \texttt{superset} check (expected
        $\subseteq$ actual) would allow extensions without penalty.

  \item \textbf{Provide a reference mock server protocol.} Understanding the
        mock's scripted response format consumed significant debugging time.
        A documented example of the request/response deck format would help.

  \item \textbf{Separate baseline tests from extension tests.} My final system
        intentionally registers many extra tools, web access, skills, memory,
        and hooks. A separate extension phase would let students demonstrate
        extra engineering work without conflicting with tests that assume the
        minimal baseline tool set.

  \item \textbf{Introduce a gradual difficulty curve for extensions.} The jump
        from ``basic agent'' to ``full coding assistant with memory, sessions,
        skills, and subagents'' is large. Scaffolding one extension as a
        guided tutorial would improve learning outcomes.
\end{enumerate}

% ──────────────────────────────────────────────────────────

\section{Conclusion}
% ──────────────────────────────────────────────────────────

The final implementation is more than a minimal ReAct loop: it is a small
runtime for terminal-based AI assistance. The core OS ideas from the course map
naturally onto the design. The tool registry resembles a syscall table, the
executor schedules read-only work in parallel while preserving ordered results,
the sandbox enforces path containment, session storage uses crash-resistant
atomic persistence, and context management behaves like a memory reclamation
layer with lossless offload followed by lossy summary.

The extensions were built to support long-running, stateful use rather than
one-shot demos. Memory makes stable user facts and preferences available across
new sessions, session directories make conversations replayable and isolate
offloaded artifacts, skills provide progressive capability loading without
inflating the base prompt, and web/browser support lets the agent inspect
external resources when local context is insufficient. These features made the
system more complex, but they also forced clearer invariants: tool results must
remain recoverable, persisted sessions must be replayable, and background
subsystems must not corrupt the main conversation.

The behavior-oriented public tests pass, including context offload and summary
(\texttt{test-ca}, \texttt{test-cb}, and C unit tests).\footnote{\footnotesize
In the aggregate \texttt{make test} run, the remaining public mismatch is the
baseline \texttt{tools\_advertised} check. It asserts that exactly the four
required tools are advertised, while this extended version deliberately exposes
sixteen tools. This is a strict baseline-interface assertion, not a failure of
offload, summary, parallel execution, or file operations.} Overall, the
project demonstrates how OS-level design principles can make an LLM agent more
bounded, recoverable, and maintainable in real interactive use.

\end{document}
